\begin{mistake}{2026-04-21}{07}

\begin{problemcontent}
设 \(A\) 为 \(m \times n\) 矩阵，\(e = [1, 1, \cdots, 1]^T\)。若方程组 \(Ay=e\) 有解，则对于 (I) \(A^T x = 0\) 与 (II) \(\begin{cases} A^T x = 0 \\ e^T x = 0 \end{cases}\) 说法正确的是（\quad）。

(A) (I) 的解都是 (II) 的解，但 (II) 的解未必是 (I) 的解 \\
(B) (II) 的解都是 (I) 的解，但 (I) 的解未必是 (II) 的解 \\
(C) (I) 的解不是 (II) 的解，且 (II) 的解也不是 (I) 的解 \\
(D) (I) 的解都是 (II) 的解，且 (II) 的解也都是 (I) 的解
\end{problemcontent}

\begin{solutioncontent}
正确答案：(D)

对方程组 (II) 而言，其解向量必须同时满足 \(A^T x = 0\) 与 \(e^T x = 0\)，故 (II) 的解自然都是 (I) 的解。

另一方面，已知 \(Ay=e\) 有解，设其特解为 \(y_0\)，则 \(Ay_0=e\)。
两边取转置得：
\[
y_0^T A^T = e^T
\]
设 \(x\) 为方程组 (I) 的任意解，即满足 \(A^T x = 0\)。考察其是否满足 (II) 的附加条件 \(e^T x = 0\)：
\[
e^T x = (y_0^T A^T) x = y_0^T (A^T x) = y_0^T \cdot 0 = 0
\]
因此，(I) 的解必定满足 \(e^T x = 0\)，即 (I) 的解也都是 (II) 的解。

综上，(I) 与 (II) 的解集互为子集，两方程组同解。
\end{solutioncontent}

\mistakeknowledge{
\begin{itemize}
 \item \textbf{核心公式/定理}：矩阵转置性质 \((AB)^T = B^T A^T\)；矩阵乘法结合律 \((AB)C = A(BC)\)。
 \item \textbf{易错点}：直觉上认为 (II) 增加了约束方程会导致解集严格变小，未能通过代入推导发现 \(e^T x=0\) 实为冗余条件。
 \item \textbf{关联知识}：从正交补空间几何意义理解，\(Ay=e\) 说明 \(e\) 属于 \(A\) 的列空间；而 \(A^T x=0\) 说明 \(x\) 与 \(A\) 的列空间正交，故 \(x\) 必然与该空间内的向量 \(e\) 正交。
\end{itemize}}

\end{mistake}
